\scanpage{46}
% archive sheet number: 45
un tel \unsure{relèvt|relèvement} existe et est unique (à isom \inserted{\unsure{(non unique|canonique|unique}} près). \deleted{\illegible{1}} Pour simplifier et pouvoir appliquer la théorie des déformations, on va supposer que $\pi A$ est stable par puissances divisées. Donc remonter $G_0$ en un gpe de B.T $G$ sur $A$ (pour le moment sans nous préoccuper de l'action de $A$) revient à remonter la filtration de $M\otimes_W k=\deleted{\unsure{M}}M_A\otimes_A k$ en une filtration de $M_A$ (sur $A$).

% Left of the following paragraph. The first row consists of three unconnected labels.
\diagram{\begin{tikzcd}[ampersand replacement=\&,column sep=small,row sep=small]
\Omega\&M=\Omega\otimes_{\mathbb Z_p}W\&M_A=\Omega\otimes_{\mathbb Z_p}A\\
A\arrow[r,dash]\arrow[d,dash]\&A\otimes_{\mathbb Z_p}W\arrow[r,dash]\arrow[d,dash]\&A\otimes_{\mathbb Z_p}A\arrow[d,dash]\\
\mathbb Z_p\arrow[r,dash]\&W\arrow[r,dash]\&A
\end{tikzcd}}

La compatibilité à l'action de $A$ signifie que c'est une filtration de $A\otimes_{\mathbb Z_p}A$-modules, i.e. on a un sous-\inserted{$A\otimes_{\mathbb Z_p}A$-}module \deleted{de rang $1$} de $\Omega\otimes_{\mathbb Z_p}A$, \inserted{de rang $1$ sur $1\otimes A$} ce dernier sera \deleted{\unsure{donc}} d'ailleurs can. isom à $\omega_G$

Or
\diagram{\begin{tikzcd}[ampersand replacement=\&,column sep=small]
\unsure{\cdot|0}\arrow[r,dash]\&\omega_G\arrow[r]\&\Omega\otimes_{\mathbb Z_p}A\arrow[r]\&t_{G^*}\arrow[r]\&0
\end{tikzcd}}
\unsure{ou|i.e.|et}, par transposition
\diagram{\begin{tikzcd}[ampersand replacement=\&,column sep=small,row sep=small]
\unsure{(***)|(**)}\&\cdot\arrow[r,dash]\&\omega_{G^*}\arrow[r,dash]\&\Omega^\vee\otimes_{\mathbb Z_p}A\arrow[r,"{\deleted{\unsure{\beta}}\rho}"]\&t_G\arrow[r,dash]\&0\\
\&\&\&A\otimes_{\mathbb Z_p}A\arrow[r,dash,"\rho_0"]\&A\&
\end{tikzcd}}

L'hypothèse que l'action de $A$ \inserted{sur $\omega_G$} (via action sur $G$) est l'action de $A$-module naturelle, signifie que \deleted{\unsure{$\rho$ est compatible|soit compatible}} \inserted{la structure} de $A\otimes_{\mathbb Z_p}A$-module sur $\omega_G$ (ou $t_G$) \deleted{\unsure{avec $\rho$ i.e. $\rho_0$ sur|avec $\rho$ i.e. $p_0$ sur}} provient de sa structure de $A$-module \deleted{triviale} évidente par restriction des scalaires à l'aide de
\[
\rho_0:A\otimes_{\mathbb Z_p}A\mathrel{\unsure{\relbar|\longrightarrow}} A\qquad \rho_0(\lambda)\otimes\mu)=\lambda\mu
\]
